\subsection{Уравнения движения механической части}

Механическая часть модуля поворота состоит из ротора электродвигателя, редуктора и выходного звена. Кинематическая схема модуля поворота представлена на рисунке~\ref{fig:kinematic_scheme}.

\begin{figure}[H]
    \centering
    \begin{tikzpicture}[
    scale=1.0,
    transform shape,
    >=Stealth,
    line/.style={draw, thick},
    arrow/.style={draw, -Stealth, thick},
    inertia/.style={
        draw,
        rectangle,
        minimum width=0.95cm,
        minimum height=0.55cm,
        align=center
    }
]

    % Макрос опоры: \support{x}{y}
    \newcommand{\support}[2]{
        % вертикальные стойки
        \draw[line] (#1-0.45,#2-0.28) -- (#1-0.45,#2+0.28);
        \draw[line] (#1+0.45,#2-0.28) -- (#1+0.45,#2+0.28);


        % верхняя штриховка
        \foreach \dx in {-0.06,0.14,0.34,0.54} {
            \draw[line]
                (#1+\dx-0.38,#2+0.55)
                -- ++(0.18,-0.18);
        }

        % верхняя направляющая
        \draw[line] (#1-0.28,#2+0.37) -- (#1+0.34,#2+0.37);

        % нижняя направляющая
        \draw[line] (#1-0.28,#2-0.37) -- (#1+0.34,#2-0.37);

        % нижняя штриховка
        \foreach \dx in {-0.06,0.14,0.34,0.54} {
            \draw[line]
                (#1+\dx-0.38,#2-0.55)
                -- ++(0.18,0.18);
        }
    }

    % Инерции
    \node[inertia] (jd) at (0,0) {$J_{\text{д}}$};
    \node[inertia] (jm) at (9.4,0) {$J_{\text{м}}$};

    % Левая часть
    \draw[line] (jd.east) -- (3.7,0);
    \support{2.6}{0}

    % Редуктор
    \draw[line] (3.7,-3) -- (3.7,0.65);
    \draw[line] (5.5,-3) -- (5.5,0.65);

    \draw[line] (3.7-0.3,0.65) -- (3.7+0.3,0.65);
    \draw[line] (3.7-0.3,-0.65) -- (3.7+0.3,-0.65);
    \draw[line] (5.5-0.3,0.65) -- (5.5+0.3,0.65);
    \draw[line] (5.5-0.3,-0.65) -- (5.5+0.3,-0.65);

    \draw[line] (3.7,-2) -- (5.5,-2);

    \support{4.6}{-2}

    \node at (4.6,0.18) {$i$};

    % Правая часть
    \draw[line] (5.5,0) -- (jm.west);
    \support{6.5}{0}

    % Момент двигателя
    \draw[arrow]
        (-1.55,-0.65)
        arc[start angle=-90,end angle=90,radius=0.65cm];

    \node at (-1.85,0.75) {$M_{\text{д}}$};

    % Угол ротора
    \draw[arrow]
        (1.05,-0.42)
        arc[start angle=-75,end angle=105,radius=0.42cm];

    \node at (1.1,0.68) {$\varphi_{\text{р}}$};

    % Угол выходного звена
    \draw[arrow]
        (8.2,-0.42)
        arc[start angle=-75,end angle=105,radius=0.42cm];

    \node at (8.25,0.68) {$\varphi_{\text{м}}$};

\end{tikzpicture}

    \caption{Кинематическая схема модуля поворота}
    \label{fig:kinematic_scheme}
\end{figure}

Для жесткой системы углы поворота ротора двигателя и выходного звена связаны соотношением
\[
    \varphi_{\text{р}} = i \varphi_{\text{м}},
\]
где $i$~--- передаточное отношение редуктора.

Введем приведенную координату ротора двигателя
\[
    \varphi_{\text{рпр}} = \frac{\varphi_{\text{р}}}{i}.
\]
Тогда для жесткой системы
\[
    \varphi_{\text{рпр}} = \varphi_{\text{м}}.
\]

Запишем кинетическую энергию системы:
\[
    T =
    \frac{1}{2} J_{\text{д}} \dot{\varphi}_{\text{р}}^2
    +
    \frac{1}{2} J_{\text{м}} \dot{\varphi}_{\text{м}}^2.
\]

Так как
\[
    \dot{\varphi}_{\text{р}} = i \dot{\varphi}_{\text{м}},
\]
то
\[
    T =
    \frac{1}{2} J_{\text{д}} i^2 \dot{\varphi}_{\text{м}}^2
    +
    \frac{1}{2} J_{\text{м}} \dot{\varphi}_{\text{м}}^2
    =
    \frac{1}{2}
    \left(
        J_{\text{дпр}} + J_{\text{м}}
    \right)
    \dot{\varphi}_{\text{м}}^2,
\]
где
\[
    J_{\text{дпр}} = J_{\text{д}} i^2
\]
--- момент инерции ротора двигателя, приведенный к координате $\varphi_{\text{м}}$.

Движущий момент электродвигателя является внешним моментом, приложенным к ротору двигателя. Элементарная работа внешних сил на возможном перемещении:
\[
    \delta A =
    M_{\text{д}} \delta \varphi_{\text{р}}
    =
    M_{\text{д}} i \delta \varphi_{\text{м}}
    =
    M_{\text{дпр}} \delta \varphi_{\text{м}},
\]
где
\[
    M_{\text{дпр}} = M_{\text{д}} i
\]
--- движущий момент двигателя, приведенный к координате $\varphi_{\text{м}}$.

Учтем упругость редуктора. Динамическая модель механической части представлена на рисунке~\ref{fig:dynamic_model}.

\begin{figure}[H]
    \centering
    \begin{tikzpicture}[
    scale=0.95,
    transform shape,
    >=Stealth,
    line/.style={draw, thick},
    arrow/.style={draw, -Stealth, thick},
    mass/.style={draw, rectangle, minimum width=0.75cm, minimum height=1.15cm},
    spring/.style={thick, decorate, decoration={zigzag, segment length=5pt, amplitude=3pt}},
    damper/.style={thick}
]

    % Массы
    \node[mass] (jdpr) at (0,0) {$J_{\text{дпр}}$};
    \node[mass] (jm) at (6.5,0) {$J_{\text{м}}$};

    % Валы
    \draw[line] (-1.0,0) -- (jdpr.west);
    \draw[line] (jdpr.east) -- (2.0,0);
    \draw[line] (4.5,0) -- (jm.west);
    \draw[line] (jm.east) -- (7.4,0);

    % Упругий и демпфирующий элементы
    \draw[line] (2.0,0.35) -- (2.45,0.35);
    \draw[spring] (2.45,0.35) -- (4.05,0.35);
    \draw[line] (4.05,0.35) -- (4.5,0.35);
    \node at (3.25,0.85) {$c$};

    \draw[line] (2.0,-0.35) -- (2.75,-0.35);
    \draw[damper] (2.75,-0.65) -- (2.75,-0.05);
    \draw[damper] (2.75,-0.65) -- (3.45,-0.65);
    \draw[damper] (2.75,-0.05) -- (3.45,-0.05);
    \draw[damper] (3.45,-0.50) -- (3.45,-0.20);
    \draw[line] (3.45,-0.35) -- (4.5,-0.35);
    \node at (3.25,-0.95) {$b$};

    % Момент двигателя
    \draw[arrow] (-1.05,-0.65) arc[start angle=-90,end angle=90,radius=0.65cm];
    \node at (-1.45,0.85) {$M_{\text{дпр}}$};

    % Углы
    \draw[arrow] (0.9,0.55) arc[start angle=30,end angle=150,radius=0.5cm];
    \node at (1.0,0.9) {$\varphi_{\text{рпр}}$};

    \draw[arrow] (5.8,0.55) arc[start angle=150,end angle=30,radius=0.5cm];
    \node at (5.9,0.9) {$\varphi_{\text{м}}$};

    % Моменты в редукторе
    \draw[arrow] (1.55,-0.55) arc[start angle=-30,end angle=-150,radius=0.45cm];
    \node at (1.7,-0.95) {$M_{\text{м}}$};

    \draw[arrow] (5.05,-0.55) arc[start angle=-150,end angle=-30,radius=0.45cm];
    \node at (5.0,-0.95) {$M_{\text{м}}$};

\end{tikzpicture}

    \caption{Динамическая модель механической части}
    \label{fig:dynamic_model}
\end{figure}

Момент на выходном валу редуктора связан с деформацией $\varphi_{\text{рпр}} - \varphi_{\text{м}}$ линейным упруго-диссипативным соотношением:
\[
    M_{\text{м}}
    =
    c \left(
        \varphi_{\text{рпр}} - \varphi_{\text{м}}
    \right)
    +
    b \left(
        \dot{\varphi}_{\text{рпр}} - \dot{\varphi}_{\text{м}}
    \right),
\]
где $c$~--- коэффициент жесткости редуктора, приведенный к выходному валу; $b$~--- коэффициент демпфирования.

Составим систему уравнений движения механической части:
\[
    \begin{cases}
        J_{\text{дпр}} \ddot{\varphi}_{\text{рпр}}
        =
        M_{\text{дпр}} - M_{\text{м}},
        \\
        J_{\text{м}} \ddot{\varphi}_{\text{м}}
        =
        M_{\text{м}}.
    \end{cases}
\]

Подставляя выражение для $M_{\text{м}}$, получим:
\[
    \begin{cases}
        J_{\text{дпр}} \ddot{\varphi}_{\text{рпр}}
        +
        b \dot{\varphi}_{\text{рпр}}
        -
        b \dot{\varphi}_{\text{м}}
        +
        c \varphi_{\text{рпр}}
        -
        c \varphi_{\text{м}}
        -
        M_{\text{дпр}}
        =
        0,
        \\
        J_{\text{м}} \ddot{\varphi}_{\text{м}}
        -
        b \dot{\varphi}_{\text{рпр}}
        +
        b \dot{\varphi}_{\text{м}}
        -
        c \varphi_{\text{рпр}}
        +
        c \varphi_{\text{м}}
        =
        0.
    \end{cases}
\]
